 % 地支相刑图
 % 用途：展示寅巳申三刑、丑未戌三刑、子卯相刑、辰午酉亥自刑
 \documentclass[tikz,border=10pt]{standalone}
 \usepackage{ctex}
 \usetikzlibrary{arrows.meta, positioning, calc, decorations.pathmorphing}
 
 \begin{document}
 \begin{tikzpicture}[scale=1.2]
   \coordinate (zi)  at ( 90:3.5);
   \coordinate (chou) at ( 60:3.5);
   \coordinate (yin)  at ( 30:3.5);
   \coordinate (mao)  at (  0:3.5);
   \coordinate (chen) at (330:3.5);
   \coordinate (si)   at (300:3.5);
   \coordinate (wu)   at (270:3.5);
   \coordinate (wei)  at (240:3.5);
   \coordinate (shen) at (210:3.5);
   \coordinate (you)  at (180:3.5);
   \coordinate (xu)   at (150:3.5);
   \coordinate (hai)  at (120:3.5);
 
   \draw[thick, gray!60] (0,0) circle (3.5);
 
   \foreach \name/\angle/\text in {
     zi/90/子, chou/60/丑, yin/30/寅, mao/0/卯,
     chen/330/辰, si/300/巳, wu/270/午, wei/240/未,
     shen/210/申, you/180/酉, xu/150/戌, hai/120/亥
   } {
     \fill[black] (\angle:3.5) circle (2pt);
     \node[font=\large] at (\angle:4.2) {\text};
   }
 
   % 寅巳申三刑：虚线三角
   \draw[brown!80!black, thick, dashed, {Stealth[length=3mm]}-{Stealth[length=3mm]}] (yin)  -- (si);
   \draw[brown!80!black, thick, dashed, {Stealth[length=3mm]}-{Stealth[length=3mm]}] (si)   -- (shen);
   \draw[brown!80!black, thick, dashed, {Stealth[length=3mm]}-{Stealth[length=3mm]}] (shen) -- (yin);
 
   % 丑未戌三刑：虚线三角
   \draw[brown!60!black, thick, dashed, {Stealth[length=3mm]}-{Stealth[length=3mm]}] (chou) -- (wei);
   \draw[brown!60!black, thick, dashed, {Stealth[length=3mm]}-{Stealth[length=3mm]}] (wei)  -- (xu);
   \draw[brown!60!black, thick, dashed, {Stealth[length=3mm]}-{Stealth[length=3mm]}] (xu)   -- (chou);
 
   % 子卯相刑：单独连线
   \draw[pink!70!black, very thick, decorate, decoration={snake, amplitude=0.5mm, segment length=4mm}, {Stealth[length=3mm]}-{Stealth[length=3mm]}] (zi) -- (mao);
 
   % 自刑：自环标记
   \foreach \angle in {330, 300, 180, 120} {
     \draw[pink!50!black, thick] (\angle:3.5) circle (8pt);
   }
 
   % 图例
   \node[font=\footnotesize, anchor=west] at (-5.5,-5.0) {\textbf{三刑}（虚线三角）：寅巳申（深）+ 丑未戌（浅）};
   \node[font=\footnotesize, anchor=west] at (-5.5,-5.5) {\textbf{相刑}（波折线）：子卯相刑};
   \node[font=\footnotesize, anchor=west] at (-5.5,-6.0) {\textbf{自刑}（自环）：辰、午、酉、亥};
   \node[font=\footnotesize, anchor=west] at (-5.5,-6.5) {刑 = 拧巴、内耗、结构性摩擦};
 \end{tikzpicture}
 \end{document}
